% ======================================================================
\begingroup\sloppy
\section{Finite terms for an arbitrary ordered charge correlator}\label{app:finite}
% ======================================================================
% Sources: limits_checked/Finite_terms_final.tex (regular terms, Groups II, III), paper/notes/fix_soft.tex
% (Step 6: integrable form of Group I Rows III, V and Group III Rows III, V, VI).

We list the finite terms, Eq.~\eqref{eq:fin_def}, in the integrable form of Section~\ref{sec:fin_integrable}, for an arbitrary ordered correlator of the charges, in the limit $\Lambda\to0$, $\vee\to\infty$. The conventions are those of Section~\ref{sec:fin_conv}: phase $e^{-i\v k\cdot(\text{unprimed}-\text{primed})}$, complex conjugate implied for Groups~I and~III, charges in the order written, principal values, $\mathcal C_{\rm f}$ of Eq.~\eqref{eq:Cf}, $P^i=\Kv^i(x'-w')$, and
\begin{align}
\mathcal L^a\equiv\big[U^{ab'}(x')-U^{ab'}(w')\big]\rho^{b'}(x')\,.
\end{align}
Row~I.III is given in Eqs.~\eqref{eq:GIIII}--\eqref{eq:Mij} and the soft remainders of Rows~I.III, I.V, III.III, III.V and~III.VI in the block~\eqref{eq:Sblock}; they are not repeated. The finite parts of Rows~I.IV and~III.I cancel each other identically, and the rows with four charges only have no finite part.

\subsection{Group I}

\paragraph{Row I.II.} The term proportional to the LO structure, left after the cancellation of the ultraviolet $\beta_0$ terms against Row~I.V (Section~\ref{sec:dg_poles}), is
\begin{align}\label{eq:app_IIa}
\frac{d^3N^{\rm I.II}_1}{d^3k}&=\frac{1}{(2\pi)^3}\,\frac{g^4N_c}{32\pi^4}\,\frac{1}{k^+}\Big\{-\frac{67}{9}+\frac{\pi^2}{3}\Big\}\notag\\&\qquad\times\int_{w,w',x,x'}e^{-ik\cdot(w-w')}\,\Kv(x'-w')\cdot\Kv(x-w)\,\mathcal L^a\big[U^{ab}(w)-U^{ab}(x)\big]\rho^b(x)\,,
\end{align}
which with its complex conjugate is Eq.~\eqref{eq:constantLO}. In the grouping in which the colour of the gluon pair of Row~I.V is evaluated at its centre of mass, the self-energy part of Row~I.II, and with it this constant, is cancelled identically by the unresolved part of Row~I.V, Eq.~\eqref{eq:dg_SEid}; in the grouping used here the compensating terms are in $N_{2a}$ and $N_{2b\,\rm rest}$ of Row~I.V. The term from the part of $\mf A^{(3)}$ with $f^{abc}\{\rho^b,\rho^c\}$ is
\begin{align}\label{eq:app_IIb}
\frac{d^3N^{\rm I.II}_{2b}}{d^3k}&=-\frac{1}{(2\pi)^3}\,\frac{g^4}{64\pi^6}\,\frac{1}{k^+}\int_{w,w',x'}e^{-ik\cdot(w-w')}\notag\\&\qquad\times\int d^2p\,d^2q\int_{x,y}e^{iq\cdot(w-y)}\,e^{ip\cdot(y-x)}\,P^i\,\frac{1}{\v q^2}\,\mathcal R^i(\v q,\v p)\notag\\
&\qquad\times\mathcal L^a\Big[f^{hbc}\,U^{ah}(w)\,\{\rho^b(x),\rho^c(y)\}-f^{abc}\,U^{bd}(x)\,U^{ce}(y)\,\{\rho^d(x),\rho^e(y)\}\Big]\,,\notag\\
\mathcal R^i(\v q,\v p)&\equiv\frac12\Big[\frac{R_0^i(\v q,\v p)}{\v p^2}-\frac{R_0^i(\v q,\v q-\v p)}{(\v q-\v p)^2}\Big]\,,\notag\\
R_0^i(\v q,\v p)&\equiv-2q^i\log\frac{\v p^2}{\v q^2}+\frac{(2\v p^2-2\v q\cdot\v p+\v q^2)\,q^i+\v q^2p^i}{(\v q-\v p)^2}\log\frac{\v q^2}{(\v q-\v p)^2}\,.
\end{align}
The same kernel multiplies both colour structures: for the structure with $U(w)$ (operator ordering $\mf A^{(1)}\mf N_2$) the row contains in addition $2q^i\log\frac{\vee}{k^+}$, which is its term in the BFKL list. $R_0^i/(\v p^2\v q^2)$ alone is not integrable at $\v p\to0$ and $\v q\to\v p$; both colour structures are antisymmetric under $\v x\leftrightarrow\v y$, i.e.\ under $\v p\to\v q-\v p$, and in the antisymmetrised kernel $\mathcal R^i$ the two points cancel each other, which is the value assigned by the rule $\int d^2p/\v p^2=0$. $\mathcal R^i/\v q^2$ is integrable at $\v p\to0$ ($\simeq\log(1/|p|)/|p|$) and at $\v q\to0$, falls like $|p|^{-3}$, and at large $|q|$ like $\log\v q^2/\v q^2$, convergent through the phase as at LO; the coordinate-space form of this kernel has not been derived.

\paragraph{Row I.V.} Variables: $\v x$ is the position of the gluon with $p^+=\xi k^+$, $\v z$ that of the gluon with $k^+-p^+$, $s\equiv\v k\cdot(\v x-\v z)$; the soft remainders are in the block~\eqref{eq:Sblock}. The regular terms are
\begin{align}\label{eq:app_V1}
\frac{d^3N^{\rm I.V}_1}{d^3k}&=\frac{1}{(2\pi)^3}\,\frac{ig^4}{32\pi^5}\,\frac{1}{k^+}\int_{x,z,y,y',x',w'}e^{-ik\cdot(z-w')}\notag\\&\qquad\times\mathcal J(s)\,P\cdot\Kv(z-x)\;\Kv(y-z)\cdot\Kv(y'-x)\notag\\
&\qquad\times\mathcal L^a\Big[f^{adc}\,U^{db}(x)\,U^{ce}(z)\,\{\rho^e(y),\rho^b(y')\}-f^{abc}\,U^{cd}(y)\,U^{be}(y')\,\{\rho^d(y),\rho^e(y')\}\Big]\,,
\end{align}
\begin{align}\label{eq:app_V2}
\frac{d^3N^{\rm I.V}_2}{d^3k}&=\frac{1}{(2\pi)^3}\,\frac{g^4}{16\pi^5}\,\frac{1}{k^+}\int_{x,y,z,x',w'}e^{-ik\cdot(z-w')}\Bigg\{P^i\,\Kv^m(z-x)\,\mathcal F^{im}_{\rm reg}\notag\\&\qquad\quad\times\Big[-f^{adc}f^{fbe}\,U^{db}(x)\,U^{ce}(z)+N_c\,U^{af}(y)\Big]\notag\\
&\qquad+N_c\,\frac{11}{6}\,e^{-is}\,P\cdot\Kv(y-x)\,\frac{\Kv(y-z)\cdot(y-x)}{(x-z)^2}\Big[U^{af}(y)-U^{af}(x)\Big]\Bigg\}\notag\\&\qquad\times\big[U^{ab'}(x')-U^{ab'}(w')\big]\rho^{b'}(x')\,\rho^f(y)\,,
\end{align}
with
\begin{align}\label{eq:app_Freg}
\mathcal F^{im}_{\rm reg}&\equiv\int_0^1d\xi\,\frac{e^{-i\xi s}\,\mathcal N^{im}(\xi)}{\xi(y-x)^2+(1-\xi)(y-z)^2}\,,\notag\\
\mathcal N^{im}(\xi)&\equiv\delta^{im}\,\xi(1-\xi)\,\frac{(y-x)^2-(y-z)^2}{(x-z)^2}+\delta^{im}\big[(1-\xi)T_1+\xi T_2\big]\notag\\&\quad+\frac{(y-x)^2}{(y-z)^2}B_2^{im}+\frac{(y-z)^2}{(y-x)^2}B_1^{im}-Q^{im}\,,\notag\\
B_1^{im}&\equiv\frac{(y-x)^i(x-z)^m}{(x-z)^2}-\frac{(y-x)^i(y-z)^m}{2(y-z)^2}\,,\notag\\
B_2^{im}&\equiv\frac{(y-z)^i(x-z)^m}{(x-z)^2}+\frac{(y-x)^m(y-z)^i}{2(y-x)^2}\,,\notag\\
T_1&\equiv\frac{(y-z)\cdot(x-z)}{(x-z)^2}+\frac{(y-x)\cdot(y-z)}{2(y-x)^2}-\frac{(y-x)^2-(y-z)^2}{2(x-z)^2}\,,\notag\\
T_2&\equiv\frac{(y-x)\cdot(x-z)}{(x-z)^2}-\frac{(y-x)\cdot(y-z)}{2(y-z)^2}-\frac{(y-x)^2-(y-z)^2}{2(x-z)^2}\,,\notag\\
Q^{im}&\equiv\frac{(y-x)^m(x-z)^i}{(x-z)^2}-\frac{(y-x)^m(y-z)^i}{2(y-z)^2}\notag\\&\quad+\frac{(y-z)^m(x-z)^i}{(x-z)^2}+\frac{(y-x)^i(y-z)^m}{2(y-x)^2}\,.
\end{align}
The colour bracket of the first line of Eq.~\eqref{eq:app_V2} is the sum of those of $N_{2a}$ and $N_{2b\,\rm rest}$; $N_{2b\,\rm UV}$ cancels against Row~I.II. The full kernel of the row, including its soft remainders, is
\begin{align}\label{eq:app_F}
\mathcal F^{im}\equiv\mathcal F^{im}_{\rm reg}-\mathcal E(s)\,\frac{B_2^{im}}{(y-z)^2}-e^{-is}\,\mathcal E(-s)\,\frac{B_1^{im}}{(y-x)^2}\,,
\end{align}
the second and third terms being the remainders with the soft gluon at $\v x$ and at $\v z$, respectively. Eqs.~\eqref{eq:app_V1} and~\eqref{eq:app_V2} are integrable by themselves (no $1/R^2$ tail); the collinear point $\v z\to\v x$ of the brace of Eq.~\eqref{eq:app_V2} is integrable with the angular integral done first, by virtue of the $\frac{11}{6}$ term, the collinear subtraction of $N_{2b\,\rm rest}$, $-\int_0^1d\xi\,[\xi(1-\xi)-2]=\frac{11}{6}$. This term is the unintegrated subtraction term whose dimensionally regularised $\v z$ integral, $-\frac{11}{6}$ in $N_{2b\,\rm UV}$, cancels the $\frac{11}{3}$ term of Row~I.II (Section~\ref{sec:dg_poles}); its $\v z$ integral alone diverges. The remaining singularity at $\v z\to\v x$ is a pure second harmonic in the angle of $\v z-\v x$, i.e.\ conditionally convergent; the angular integral first, a rotation-invariant cut-off, is the regularisation in which the $\frac{11}{6}$ subtraction was obtained (paragraph ``Convergence'' of Section~\ref{sec:fin_integrable}).

\subsection{Group II}

With $\mathcal T^{ij}(x,z,w)\equiv\frac{\delta^{ij}}{2(z-w)^2}+\Kv^i(z-w)\Kv^j(x-z)-\Kv^i(x-w)\Kv^j(x-z)-\Kv^i(x-w)\Kv^j(z-w)$:

\paragraph{Row II.I, first part.}
\begin{align}\label{eq:app_III1a}
\frac{d^3N^{\rm II.I}_{1a}}{d^3k}&=-\frac{1}{(2\pi)^3}\,\frac{ig^4}{16\pi^5}\,\frac{1}{k^+}\int_{x,y,z,w,x',w'}e^{-ik\cdot(w-w')}\notag\\&\qquad\times P^i\,\Kv^j(y-z)\,\log\frac{(x-z)^2}{(x-w)^2}\,\mathcal T^{ij}(x,z,w)\notag\\
&\qquad\times\Big[U^{ac'}(w')\{\rho^{c'}(x'),\rho^b(y)\}-U^{c'b}(z)\,U^{ad'}(x')\,U^{c'e'}(y)\{\rho^{d'}(x'),\rho^{e'}(y)\}\Big]\notag\\
&\qquad\times\Big[f^{cbd}\,U^{ac}(w)\,\rho^d(x)-f^{acd}\,U^{cb}(z)\,U^{de}(x)\,\rho^e(x)\Big]\,,\\
\label{eq:app_III1b}
\frac{d^3N^{\rm II.I}_{1b}}{d^3k}&=+\frac{1}{(2\pi)^3}\,\frac{ig^4}{16\pi^5}\,\frac{1}{k^+}\int_{x,y,z,w,x',w'}e^{-ik\cdot(w-w')}\notag\\&\qquad\times\Kv^i(x'-w)\,\Kv^j(y-z)\,\log\frac{(x-z)^2}{(x-w')^2}\,\mathcal T^{ij}(x,z,w')\notag\\
&\qquad\times\Big[f^{c'bd'}\,U^{ac'}(w')\,\rho^{d'}(x)-f^{ac'd'}\,U^{c'b}(z)\,U^{d'e'}(x)\,\rho^{e'}(x)\Big]\notag\\
&\qquad\times\Big[U^{ac}(w)\{\rho^c(x'),\rho^b(y)\}-U^{cb}(z)\,U^{ad}(x')\,U^{ce}(y)\{\rho^d(x'),\rho^e(y)\}\Big]\,.
\end{align}
Each is integrable by itself: logarithmically at $\v x\to\v z$ and $\v x\to\v w$, like $1/|\v z-\v w|$ at $\v z\to\v w$ (the logarithm vanishes linearly), like $1/|\v y-\v z|$ at $\v y\to\v z$, and like $\log\v z^2/|\v z|^3$ at large $z$.

\paragraph{Row II.I, second part.} With $a\equiv x-w$, $a'\equiv x'-w'$, $b\equiv x-z$, $b'\equiv x'-z$, $c\equiv z-w$, $c'\equiv z-w'$, $A\equiv a^2$, etc., $\mathcal D\equiv A'B-AB'$, $\lambda\equiv\log\frac{B}{A}$, $\lambda'\equiv\log\frac{B'}{A'}$,
\begin{align}\label{eq:app_III2}
\frac{d^3N^{\rm II.I}_2}{d^3k}&=\frac{1}{(2\pi)^3}\,\frac{g^4}{8\pi^5}\,\frac{1}{k^+}\int_{z,w,w',x,x'}e^{-ik\cdot(w-w')}\,\mathfrak K_2\notag\\&\qquad\times\Big[f^{c'bd'}f^{cbd}\,U^{ac'}(w')\,U^{ac}(w)\,\rho^{d'}(x')\rho^d(x)\notag\\
&\qquad-f^{c'bd'}f^{acd}\,U^{ac'}(w')\,U^{cb}(z)\,U^{de}(x)\,\rho^{d'}(x')\rho^e(x)\notag\\
&\qquad-f^{ac'd'}f^{cbd}\,U^{ac}(w)\,U^{c'b}(z)\,U^{d'e'}(x')\,\rho^{e'}(x')\rho^d(x)\notag\\&\qquad+N_c\,U^{de'}(x')\,U^{de}(x)\,\rho^{e'}(x')\rho^e(x)\Big]\,,
\end{align}
$\mathfrak K_2\equiv\mathfrak K_{2a}+\mathfrak K_{2b}+\mathfrak K_{2c}+\mathfrak K_{2d}+\mathfrak K_{2e}+\mathfrak K_{2f}$, with
\begin{align}\label{eq:app_K2}
\mathfrak K_{2a}&=\frac{\lambda-\lambda'}{\mathcal D}\Bigg[\frac{a\cdot c'\,b'\cdot c}{CC'}-\frac{a\cdot a'\,c\cdot b'}{2CA'}+\frac{b\cdot b'\,a\cdot c'}{2BC'}-\frac{b\cdot b'\,a\cdot a'}{4BA'}\notag\\&\qquad\qquad+\frac{a'\cdot c\,b\cdot c'}{CC'}+\frac{a'\cdot c\,b\cdot b'}{2CB'}-\frac{b\cdot c'\,a\cdot a'}{2AC'}-\frac{a\cdot a'\,b\cdot b'}{4AB'}\Bigg]\,,\notag\\
\mathfrak K_{2b}&=\frac{1}{\mathcal D}\Big[\frac{B'}{A'}\lambda'-\frac{B}{A}\lambda\Big]\Bigg[\frac{a\cdot a'\,c\cdot c'}{CC'}+\frac{a\cdot a'\,b'\cdot c}{2CB'}\notag\\&\qquad\qquad+\frac{a\cdot a'\,b\cdot c'}{2BC'}+\frac{b\cdot b'\,a\cdot a'}{4BB'}\Bigg]\,,\notag\\
\mathfrak K_{2c}&=\frac{1}{\mathcal D}\Big[\frac{A'}{B'}\lambda'-\frac{A}{B}\lambda\Big]\Bigg[\frac{b\cdot b'\,c\cdot c'}{CC'}-\frac{b\cdot b'\,a'\cdot c}{2CA'}\notag\\&\qquad\qquad-\frac{a\cdot c'\,b\cdot b'}{2AC'}+\frac{b\cdot b'\,a\cdot a'}{4AA'}\Bigg]\,,\notag\\
\mathfrak K_{2d}&=\frac{1}{\mathcal D}\Big[\frac{B\lambda}{B-A}+\frac{B'\lambda'}{A'-B'}\Big]\Bigg[\frac{a\cdot c\,(B'-A')}{2CC'}+\frac{a\cdot b\,(B'-A')}{4BC'}\notag\\&\qquad\qquad+\frac{a'\cdot c'\,(B-A)}{2CC'}+\frac{a'\cdot b'\,(B-A)}{4CB'}\Bigg]\,,\notag\\
\mathfrak K_{2e}&=-\frac{1}{\mathcal D}\Big[\frac{A\lambda}{B-A}+\frac{A'\lambda'}{A'-B'}\Big]\Bigg[\frac{b\cdot c\,(B'-A')}{2CC'}-\frac{a\cdot b\,(B'-A')}{4AC'}\notag\\&\qquad\qquad+\frac{b'\cdot c'\,(B-A)}{2CC'}-\frac{a'\cdot b'\,(B-A)}{4CA'}\Bigg]\,,\notag\\
\mathfrak K_{2f}&=\frac{1}{2CC'}\Bigg[-\frac{AB(B'-A')}{(B-A)\,\mathcal D}\,\lambda-1+\frac{A'B'(B-A)}{(B'-A')\,\mathcal D}\,\lambda'\Bigg]\,.
\end{align}
The $-1$ in $\mathfrak K_{2f}$ is the power term $k^+[\frac1\vee-\frac{1}{k^+}]\to-1$. The two terms of each of $\mathfrak K_{2b},\dots,\mathfrak K_{2f}$ have to stay together: on the surface $\mathcal D=0$ each has a simple pole, and the residues cancel within each $\mathfrak K$ (in $\mathfrak K_{2a}$ the factor $\lambda-\lambda'$ vanishes there); on $B=A$ and $B'=A'$ the logarithms vanish. The three terms of $\mathfrak K_{2f}$ cancel at $\v z\to\v w$ and $\v z\to\v w'$. All six kernels have to be integrated over $\v z$ together: single kernels fall like $1/\v z^2$ or $\log\v z^2/\v z^2$, their sum like $\log\v z^2/|\v z|^3$. The remaining singularities are integrable.

\paragraph{Row II.IV.} The DGLAP and running-coupling terms are Eqs.~\eqref{eq:DGLAP} and~\eqref{eq:RC}, the remnant Eq.~\eqref{eq:remnant} (coordinate-space form), or, equivalently, the local terms~\eqref{eq:dg_local} and the remainder~\eqref{eq:dg_R} (local form; relation Eq.~\eqref{eq:dg_equiv}). The finite term is
\begin{align}\label{eq:app_IVfin}
\frac{d^3N^{\rm II.IV}_{\rm fin}}{d^3k}&=2N_c\,\mathcal C\int_0^1d\xi\,\frac{1-\xi}{\xi}\int_{x,y,x',y'}e^{-ik\cdot(y-y')\frac1\xi}\notag\\&\qquad\times\Bigg[\Kv(x-y)\cdot\Kv(x'-y')-2\,\frac{\Kv(x-y)\cdot(y'-y)\;\Kv(x'-y')\cdot(y'-y)}{(y'-y)^2}\Bigg]\notag\\
&\qquad\times\big[U^{ce'}(y')-U^{ce'}(x')\big]\big[U^{ce}(y)-U^{ce}(x)\big]\rho^{e'}(x')\rho^e(x)\,.
\end{align}
After the $\v y-\v y'$ integral its integrand is $O(\xi)$ at $\xi\to0$; at fixed coordinates the $\xi$ integral converges conditionally, $\int_0^1d\xi\,\frac{1-\xi}{\xi}\,e^{-ia/\xi}=(1+ia)E_1(ia)-e^{-ia}$ with $a=\v k\cdot(\v y-\v y')$, with an integrable logarithm on the line $\v k\cdot(\v y-\v y')=0$, and gives the same result as the order with the $\v y-\v y'$ integral first (the bound $|(1+ia)E_1(ia)-e^{-ia}|\le C(1+|\log|a||)$ is uniform in a lower cut-off on $\xi$; on a model with the structure of the term the two orders agree to $5\times10^{-11}$). The tensor term contributes the part $\propto\cos2\varphi$ of the tail~\eqref{eq:dg_tailtot}.

\subsection{Group III}

Colour structures (the index $a$ contracted with $f^{c'd'a}$ in Rows~III.III and~III.VI):
\begin{align}\label{eq:app_G3col}
\mathcal L^{c'd'b}(y')&\equiv U^{c'e'}(y')\,U^{d'b}(z)\,\rho^{e'}(x')-U^{d'b}(z)\,U^{c'e'}(x')\,\rho^{e'}(x')\,,\notag\\
\mathcal L^{ab}_2(y')&\equiv f^{c'd'a}\mathcal L^{c'd'b}(y')\,,\notag\\
\mathcal R_1^{ab}&\equiv U^{ac}(w)\{\rho^c(y),\rho^b(x)\}-U^{cb}(z)\,U^{ad}(y)\,U^{ce}(x)\{\rho^d(y),\rho^e(x)\}\,,\notag\\
\mathcal R_2^{ab}&\equiv f^{cbd}\,U^{ac}(w)\,\rho^d(x)-f^{acd}\,U^{cb}(z)\,U^{de}(x)\,\rho^e(x)\,,\notag\\
\mathcal L_{\rm V}^{c'}(y')&\equiv U^{c'e'}(y')\,\rho^{e'}(x')-U^{c'e'}(x')\,\rho^{e'}(x')\,,\notag\\
\mathcal R_{\rm V}^{ac}&\equiv U^{ce}(x)\,U^{ad}(y)\,\rho^e(x)\rho^d(y)-U^{ad}(w)\,U^{ce}(x)\,\rho^e(x)\rho^d(y)\,,\notag\\
\mathcal R_{\rm VI}^{ab}&\equiv U^{cb}(z)\,U^{ad}(y)\,U^{ce}(x)\,\rho^d(y)\,\rho^e(x)-U^{ad}(y)\,\rho^d(y)\,\rho^b(x)\,;
\end{align}
$\mathcal L(z)$ denotes $\mathcal L(y')$ with $U^{c'e'}(y')\to U^{c'e'}(z)$. In Rows~III.III, III.V and~III.VI, $s\equiv\v k\cdot(\v z-\v y')$; the soft remainders ($\Psi$ terms) are in the block~\eqref{eq:Sblock}.

\paragraph{Row III.II.}
\begin{align}\label{eq:app_IIIII}
\frac{d^3N^{\rm III.II}}{d^3k}&=-\frac{1}{(2\pi)^3}\,\frac{ig^4}{8\pi^5}\,\frac{1}{k^+}\int_{x,y,x',z,w,w'}e^{-ik\cdot(w-w')}\notag\\&\qquad\times P^i\,\Kv^j(y-z)\,\log\frac{(x-z)^2}{(x-w)^2}\,\mathcal T^{ij}(x,z,w)\notag\\
&\qquad\times\Big[U^{c'd'}(y)\,U^{ab'}(x')\,U^{c'b}(z)\,\rho^{d'}(y)\rho^{b'}(x')-U^{ab'}(x')\,\rho^b(y)\rho^{b'}(x')\Big]\notag\\
&\qquad\times\Big[f^{cbd}\,U^{ac}(w)\,\rho^d(x)-f^{acd}\,U^{cb}(z)\,U^{de}(x)\,\rho^e(x)\Big]\,.
\end{align}
Its kernel is that of $N_{1a}$ of Row~II.I, with the same integrability.

\paragraph{Row III.III, first part.}
\begin{align}\label{eq:app_III31}
\frac{d^3N^{\rm III.III}_1}{d^3k}&=\frac{1}{(2\pi)^3}\,\frac{ig^4f^{c'd'a}}{16\pi^5}\,\frac{1}{k^+}\int_{w,x',y',z,y,x}e^{-ik\cdot(w-y')}\notag\\
&\qquad\times\Bigg\{\widehat{\mathcal E}(s)\Big[\Kv(y'-z)\cdot\Kv(x'-y')\,\Kv(y-w)\cdot\Kv(x-z)\notag\\&\qquad\qquad-K(y',x,z)\,\Kv(x'-y')\cdot\Kv(y-w)\Big]\mathcal L^{c'd'b}(y')\notag\\
&\qquad-\Big[\pi\delta(s)-i\mathcal P\frac1s\Big]\Kv(y'-z)\cdot\Kv(y-w)\,\Kv^j(x-z)\notag\\
&\qquad\qquad\times\Big[\Kv^j(x'-y')\,\mathcal L^{c'd'b}(y')-\Kv^j(x'-z)\,\mathcal L^{c'd'b}(z)\Big]\Bigg\}\mathcal R_1^{ab}\,.
\end{align}

\paragraph{Row III.III, second part.} With $a\equiv x-w$, $b\equiv x-z$, $c\equiv z-w$, $A\equiv a^2$, $B\equiv b^2$, $C\equiv c^2$, $D_p\equiv p^+B+k^+A$,
\begin{align}\label{eq:app_IJW}
\mathcal I&\equiv\frac{a\cdot\Kv(x'-y')\;c\cdot\Kv(y'-z)}{C}+\frac{b\cdot\Kv(y'-z)\;a\cdot\Kv(x'-y')}{2B}\,,\notag\\
\mathcal W^{k'm}&\equiv\frac{b^{k'}c^m}{C}-\frac{b^{k'}a^m}{2A}\,,\notag\\
\mathcal V&\equiv\frac{b\cdot\Kv(y'-z)\;c\cdot\Kv(x'-y')}{C}-\frac{b\cdot\Kv(y'-z)\;a\cdot\Kv(x'-y')}{2A}\notag\\&\quad+\frac{a\cdot\Kv(y'-z)\;c\cdot\Kv(x'-y')}{C}+\frac{b\cdot\Kv(x'-y')\;a\cdot\Kv(y'-z)}{2B}\,,
\end{align}
\begin{align}\label{eq:app_III32}
\frac{d^3N^{\rm III.III}_2}{d^3k}&=\frac{1}{(2\pi)^3}\,\frac{g^4}{4\pi^4}\int_0^\infty\frac{dp^+}{2\pi}\int_{w,x',y',z,x}e^{-ik\cdot(w-y')}\notag\\&\qquad\qquad\times e^{-i\frac{p^+}{k^+}s}\,\frac{1}{D_p}\,\{\dots\}_{\rm reg}\,\mathcal L_2^{ab}(y')\,\mathcal R_2^{ab}\notag\\
&\quad-\frac{1}{(2\pi)^3}\,\frac{g^4}{4\pi^4}\,\frac{1}{k^{+2}}\int_0^\infty\frac{dp^+}{2\pi}\int_{w,x',y',z,x}e^{-ik\cdot(w-y')}\notag\\&\qquad\qquad\times e^{-i\frac{p^+}{k^+}s}\,\frac{p^+}{D_p}\,\Kv^m(y'-z)\,\mathcal W^{k'm}\notag\\
&\qquad\qquad\times\Big[\Kv^{k'}(x'-y')\,\mathcal L_2^{ab}(y')-\Kv^{k'}(x'-z)\,\mathcal L_2^{ab}(z)\Big]\mathcal R_2^{ab}\notag\\
&\quad+\frac{1}{(2\pi)^3}\,\frac{g^4}{8\pi^5}\,\frac{1}{k^+}\int_{w,x',w',z,x}e^{-ik\cdot(w-w')}\notag\\&\qquad\qquad\times\frac{(x-w)^2}{(x-z)^2[(x-w)^2-(x-z)^2]}\log\frac{(x-w)^2}{(x-z)^2}\notag\\&\qquad\qquad\times\Kv^m(w'-z)\,\Kv^{k'}(x'-z)\,\mathcal W^{k'm}\,\mathcal L_2^{ab}(z)\,\mathcal R_2^{ab}\,,\notag\\
\{\dots\}_{\rm reg}&\equiv\Kv(y'-z)\cdot\Kv(x'-y')\notag\\&\qquad\times\bigg[\frac{p^+(B-A)}{C\,(p^++k^+)^2}\notag\\&\qquad+\frac{p^+}{k^+(p^++k^+)}\Big(\frac{b\cdot c}{C}-\frac{a\cdot b}{2A}-\frac{B-A}{2C}\Big)\notag\\
&\qquad+\frac{1}{p^++k^+}\Big(\frac{a\cdot c}{C}+\frac{a\cdot b}{2B}-\frac{B-A}{2C}\Big)\bigg]\notag\\&\quad+\frac{B-A}{(p^++k^+)A}\,\mathcal I-\frac{1}{k^+}\,\mathcal V\,.
\end{align}
In the second term the $\v y'$ integral is done before the $p^+$ integral. The coefficient of $\mathcal I$ in the brace is the plus-prescription term with the phase kept in the subtraction plus the term $\frac{1}{p^++k^+}$, $-\big[\frac{k^+}{p^+(p^++k^+)}(1-\frac{D_p}{k^+A})+\frac{1}{p^++k^+}\big]=\frac{B-A}{(p^++k^+)A}$, so that no $1/p^+$ is left; the brace has to stay under one $p^+$ and one $\v z$ integral (at large $z$ and small $p^+$ the term with $\mathcal I$ and one term of $\mathcal V$ are separately $\propto1/(p^+\v z^2+k^+A)$, and only their sum is integrable). The last term is what is left of the term added back for the subtraction in the second term after its $\frac{1}{(x-z)^2}\log\frac{\vee}{k^+}$ has been assigned to the BFKL term. The first term falls like $\log R/R^3$ at large $z$.

\paragraph{Rows III.V and III.VI.}
\begin{align}\label{eq:app_III5}
\frac{d^3N^{\rm III.V}}{d^3k}&=\frac{1}{(2\pi)^3}\,\frac{ig^4f^{c'ca}}{8\pi^5}\,\frac{1}{k^+}\int_{w,y',z,y,x,x'}e^{-ik\cdot(w-y')}\notag\\
&\qquad\times\Bigg\{\widehat{\mathcal E}(s)\Big[\Kv(y'-z)\cdot\Kv(x'-y')\,\Kv(y-w)\cdot\Kv(x-z)\notag\\&\qquad\qquad-K(y',x,z)\,\Kv(x'-y')\cdot\Kv(y-w)\Big]\mathcal L_{\rm V}^{c'}(y')\notag\\
&\qquad-\Big[\pi\delta(s)-i\mathcal P\frac1s\Big]\Kv(y'-z)\cdot\Kv(y-w)\,\Kv^j(x-z)\notag\\
&\qquad\qquad\times\Big[\Kv^j(x'-y')\,\mathcal L_{\rm V}^{c'}(y')-\Kv^j(x'-z)\,\mathcal L_{\rm V}^{c'}(z)\Big]\Bigg\}\mathcal R_{\rm V}^{ac}\,.
\end{align}
Row~III.VI is Eq.~\eqref{eq:app_III5} with $f^{c'ca}\mathcal L_{\rm V}^{c'}\mathcal R_{\rm V}^{ac}\to f^{c'd'a}\mathcal L^{c'd'b}\mathcal R_{\rm VI}^{ab}$, with the same prefactor. In Rows~III.III (first part), III.V and~III.VI, $\widehat{\mathcal E}(s)\simeq\frac{1}{is}$ at large $z$ gives a $|\v z|^{-3}$ decay, with a logarithm on the line $\v k\cdot(\v z-\v y')=0$; the distributional terms are finite (Section~\ref{sec:fin_inst}) and have no tail without the phase $e^{-i\v k\cdot(\v w-\v z)}$.

\subsection{The out-state term \texorpdfstring{$\Delta N$}{Delta N} in momentum space}\label{app:DNmom}

The on-shell part of $\mf B_3$, Eq.~\eqref{eq:B3mom}, is simple in momentum space. Writing its pole part as $ig^2f^{abc}\rho^c(\v k-\v p)\,n_{ji}(k,p)/D$, $D=p^+\v k^2-k^+\v p^2$, on the energy shell $p^+=p_0=k^+\v p^2/\v k^2$ the vertex is a rotation of the polarisation by twice the scattering angle,
\begin{align}\label{eq:app_nshell}
n_{ji}(k,p)\big|_{p^+=p_0}=\frac{|\v k|\,|\v p|}{(\v k-\v p)^2}\,R(2\theta)_{ji}\,,\qquad R(\alpha)=\begin{pmatrix}\cos\alpha&-\sin\alpha\\ \sin\alpha&\cos\alpha\end{pmatrix},
\end{align}
$\theta$ the angle from $\v k$ to $\v p$ (verified numerically at $10^4$ random points, to $10^{-8}$; there is no pole on the circle $|\v p|=|\v k|$), and $n_{ji}(k,p)=n_{ij}(p,k)$ on the shell, which makes the on-shell part anti-Hermitian. In the momentum-space conventions of Li and Skokov~\cite{Li:2021ntm} (their Fourier sign, $U_{\rm LS}=U^T$), in which the LO amplitude is $\mathcal A^c_{{\rm LO}\,j}=-\frac{\sqrt2g}{\sqrt{k^+}}M^c_j(\v k)$, $M^c_j(\v k)=\int\frac{d^2p}{(2\pi)^2}L^j(\v p,\v k)\,[\rho(\v p)U_{\rm LS}(\v k-\v p)]^c$, $L(\v p,\v k)=\frac{\v p}{\v p^2}-\frac{\v k}{\v k^2}$, and $k^+\frac{d^3N_{\rm LO}}{d^3k}=\frac{g^2}{4\pi^3}|M|^2$,
\begin{align}\label{eq:app_DNmom}
k^+\frac{d^3\Delta N}{d^3k}=-\frac{\eta\,g^4}{4\pi^3}\int\frac{d^2p}{(2\pi)^2}\,\frac{R(2\theta)_{ji}}{(\v k-\v p)^2}\,{\rm Re}\Big[M^{a*}_i(\v k)\,f^{abc}\,\bar\rho^c(\v k-\v p)\,M^b_j(\v p)\Big],
\end{align}
with $\bar\rho^c(\v q)=\int\frac{d^2p}{(2\pi)^2}[\rho(\v p)U_{\rm LS}(\v q-\v p)]^c$ the charge rotated by the target. $\Delta N$ is the on-shell (Coulomb) rescattering of the produced LO gluon on the rotated charge, with the Coulomb factor $1/(\v k-\v p)^2$ and a helicity rotation. At $\v p\to\v k$ the leading term $f^{abc}{\rm Re}(M^{a*}_iM^b_i)\bar\rho^c(0)$ vanishes identically (the forward Coulomb phase), so that the integrand is $O(1/|\v k-\v p|)$ and integrable. Before this cancellation, i.e.\ for the term of Row~I.III without its complex conjugate and with the window $|p^+-k^+|>\Lambda$, the forward region gives a Coulomb logarithm $\propto\eta\log\frac{k^+}{\Lambda}\,f^{abc}\bar Q^c\mathcal A^{a*}_{{\rm LO}\,i}\mathcal A^b_{{\rm LO}\,i}$ ($\bar Q=\bar\rho(0)$), whose coefficient agrees exactly with the one that follows from the coordinate form, Eq.~\eqref{eq:DeltaN}. The full Fourier identity between Eq.~\eqref{eq:app_DNmom} and the coordinate form~\eqref{eq:DeltaN} was checked through this forward coefficient and the tensor structure (both are rotation--dilation matrices of the same homogeneity), not numerically in full.

\subsection{Cubic rows: symmetrised three-\texorpdfstring{$\rho$}{rho} part and commutator terms}\label{app:fin_sym}

Each cubic colour operator of the rows above is a $c$-number tensor times three ordered charges. Before any average over the projectile, and without any assumption on it, the ordered product is rewritten exactly as the totally symmetrised product plus single commutators, which are two-$\rho$ operators, plus double commutators, which are one-$\rho$ operators. Here $\{\dots\}_{\rm sym}$ is the average over all orderings of the operators inside the braces, the commutator is that of Eq.~\eqref{eq:rho}, $[\rho^a(\v x),\rho^b(\v y)]=if^{abc}\rho^c(\v x)\delta^{(2)}(\v x-\v y)$, and the indices $r$, $s$ are summed. For a tensor $C^{abc}$ and points $p_1,p_2,p_3$, Eq.~\eqref{eq:symmetrise} gives, for the three types that occur,
\begin{align}\label{eq:app_symord}
C^{abc}\rho^a(p_1)\rho^b(p_2)\rho^c(p_3)&=C^{abc}\{\rho^a(p_1)\rho^b(p_2)\rho^c(p_3)\}_{\rm sym}\notag\\
&\quad+\frac i2\,C^{abc}\Big[f^{abr}\,\delta^{(2)}(p_1-p_2)\,\{\rho^r(p_1)\rho^c(p_3)\}_{\rm sym}\notag\\
&\qquad+f^{bcr}\,\delta^{(2)}(p_2-p_3)\,\{\rho^a(p_1)\rho^r(p_2)\}_{\rm sym}\notag\\
&\qquad+f^{acr}\,\delta^{(2)}(p_1-p_3)\,\{\rho^r(p_1)\rho^b(p_2)\}_{\rm sym}\Big]\notag\\
&\quad-\frac16\,C^{abc}\big[f^{abr}f^{rcs}-f^{bcr}f^{ras}\big]\notag\\
&\qquad\times\delta^{(2)}(p_1-p_2)\,\delta^{(2)}(p_1-p_3)\,\rho^s(p_1)\,,\\
\label{eq:app_syma23}
C^{abc}\rho^a(p_1)\{\rho^b(p_2),\rho^c(p_3)\}&=2\,C^{abc}\{\rho^a(p_1)\rho^b(p_2)\rho^c(p_3)\}_{\rm sym}\notag\\
&\quad+i\,C^{abc}\Big[f^{abr}\,\delta^{(2)}(p_1-p_2)\,\{\rho^r(p_1)\rho^c(p_3)\}_{\rm sym}\notag\\
&\qquad+f^{acr}\,\delta^{(2)}(p_1-p_3)\,\{\rho^r(p_1)\rho^b(p_2)\}_{\rm sym}\Big]\notag\\
&\quad-\frac16\,C^{abc}\big[f^{abr}f^{rcs}+f^{acr}f^{rbs}\big]\notag\\
&\qquad\times\delta^{(2)}(p_1-p_2)\,\delta^{(2)}(p_1-p_3)\,\rho^s(p_1)\,,\\
\label{eq:app_syma12}
C^{abc}\{\rho^a(p_1),\rho^b(p_2)\}\rho^c(p_3)&=2\,C^{abc}\{\rho^a(p_1)\rho^b(p_2)\rho^c(p_3)\}_{\rm sym}\notag\\
&\quad+i\,C^{abc}\Big[f^{bcr}\,\delta^{(2)}(p_2-p_3)\,\{\rho^a(p_1)\rho^r(p_2)\}_{\rm sym}\notag\\
&\qquad+f^{acr}\,\delta^{(2)}(p_1-p_3)\,\{\rho^r(p_1)\rho^b(p_2)\}_{\rm sym}\Big]\notag\\
&\quad+\frac16\,C^{abc}\big[f^{bcr}f^{ras}+f^{acr}f^{rbs}\big]\notag\\
&\qquad\times\delta^{(2)}(p_1-p_2)\,\delta^{(2)}(p_1-p_3)\,\rho^s(p_1)\,.
\end{align}
Under (A1) and (A2) of Section~\ref{sec:fin_singlet}, $\langle\{\rho^a\rho^b\rho^c\}_{\rm sym}\rangle=S^{abc}$, $\langle\{\rho^a\rho^b\}_{\rm sym}\rangle=\delta^{ab}n/(N_c^2-1)$ and $\langle\rho^s\rangle=0$, which gives Eq.~\eqref{eq:cubicavg}: the first line of each identity yields the $(3\rho)$ equations of Section~\ref{sec:fin_singlet}, the bracket with the $\delta$-functions of two points the $(2\rho)$ (contact) equations, and the last line does not contribute. Below, each cubic row is written in this form with its own tensor. The indices $a,b,c$ are those of the charges in the order written, all other indices are summed, $\mathcal L^a=[U^{ab'}(x')-U^{ab'}(w')]\rho^{b'}(x')$, and the Group~III structures are those of Eq.~\eqref{eq:app_G3col}; in Rows~III.III and~III.VI $U(z)U^T(z)=1$ has been used. At a $\delta$-function the Wilson lines of the tensor are taken at the coincident point. The kernel and prefactor that multiply each colour operator are those of the row equation named in the heading; they are not affected by the decomposition. The decompositions as printed were checked at the operator level: with $SU(3)$ charges represented by generators on three lattice sites, the $\delta$-functions as Kronecker deltas of the sites and random Wilson lines (independent, and equal at coinciding sites), for all $27$ assignments of the three charge points of each of the thirteen colour structures to the sites, the right-hand sides reproduce the ordered products of the colour structures of the rows, read from the source of the row equations, to $3\times10^{-15}$ ($702$ cases; a single sign error gives deviations of order one).

\paragraph{Row I.II, $N_{2b}$, Eq.~\eqref{eq:app_IIb}.} The colour operator is $\mathcal L^a\big[f^{hbc}\,U^{ah}(w)\,\{\rho^b(x),\rho^c(y)\}-f^{abc}\,U^{bd}(x)\,U^{ce}(y)\,\{\rho^d(x),\rho^e(y)\}\big]=\mathcal C^{abc}_{2b}\,\rho^a(x')\{\rho^b(x),\rho^c(y)\}$, type $\rho\{\rho,\rho\}$, $(p_1,p_2,p_3)=(x',x,y)$, with
\begin{align}
\mathcal C^{abc}_{2b}&\equiv \big[U^{ga}(x')-U^{ga}(w')\big]\big[f^{hbc}\,U^{gh}(w)-f^{gde}\,U^{db}(x)\,U^{ec}(y)\big]\,,\notag\\
\mathcal C^{abc}_{2b}\,\rho^a(x')\{\rho^b(x),\rho^c(y)\}
&=2\,\mathcal C^{abc}_{2b}\,\{\rho^a(x')\rho^b(x)\rho^c(y)\}_{\rm sym}\notag\\
&\quad+i\,\mathcal C^{abc}_{2b}\,\Big[f^{abr}\,\delta^{(2)}(x'-x)\,\{\rho^{r}(x)\rho^{c}(y)\}_{\rm sym}\notag\\
&\qquad+f^{acr}\,\delta^{(2)}(x'-y)\,\{\rho^{r}(y)\rho^{b}(x)\}_{\rm sym}\Big]\notag\\
&\quad-\frac16\,\mathcal C^{abc}_{2b}\big[f^{abr}f^{rcs}+f^{acr}f^{rbs}\big]\notag\\
&\qquad\times\delta^{(2)}(x'-x)\,\delta^{(2)}(x'-y)\,\rho^s(x)\,.\label{eq:app_sym_IIb}
\end{align}
\paragraph{Row I.III, structures $C_D$ ($u=z$) and $C_A$ ($u=w$), Eq.~\eqref{eq:GIIII}.} The colour operators are $C_D,C_A=f^{abd}\,\mathcal L^a\,U^{de}(x)\big[U^{bc}(u)-U^{bc}(y)\big]\rho^e(x)\rho^c(y)=\mathcal C^{abc}_{\rm I.III}(u)\,\rho^a(x')\rho^b(x)\rho^c(y)$, ordered, $(p_1,p_2,p_3)=(x',x,y)$, with
\begin{align}
\mathcal C^{abc}_{\rm I.III}(u)&\equiv \big[U^{ga}(x')-U^{ga}(w')\big]\,f^{ghd}\,U^{db}(x)\big[U^{hc}(u)-U^{hc}(y)\big]\,,\notag\\
\mathcal C^{abc}_{\rm I.III}(u)\,\rho^a(x')\rho^b(x)\rho^c(y)
&=\mathcal C^{abc}_{\rm I.III}(u)\,\{\rho^a(x')\rho^b(x)\rho^c(y)\}_{\rm sym}\notag\\
&\quad+\frac i2\,\mathcal C^{abc}_{\rm I.III}(u)\,\Big[f^{abr}\,\delta^{(2)}(x'-x)\,\{\rho^{r}(x)\rho^{c}(y)\}_{\rm sym}\notag\\
&\qquad+f^{bcr}\,\delta^{(2)}(x-y)\,\{\rho^{a}(x')\rho^{r}(x)\}_{\rm sym}\notag\\
&\qquad+f^{acr}\,\delta^{(2)}(x'-y)\,\{\rho^{r}(y)\rho^{b}(x)\}_{\rm sym}\Big]\notag\\
&\quad-\frac16\,\mathcal C^{abc}_{\rm I.III}(u)\big[f^{abr}f^{rcs}-f^{bcr}f^{ras}\big]\notag\\
&\qquad\times\delta^{(2)}(x'-x)\,\delta^{(2)}(x'-y)\,\rho^s(x)\,.\label{eq:app_sym_IIII}
\end{align}
\paragraph{Row I.V, first part, Eq.~\eqref{eq:app_V1}.} The colour operator is $\mathcal L^a\big[f^{adc}\,U^{db}(x)\,U^{ce}(z)\,\{\rho^e(y),\rho^b(y')\}-f^{abc}\,U^{cd}(y)\,U^{be}(y')\,\{\rho^d(y),\rho^e(y')\}\big]=\mathcal C^{abc}_{\rm I.V}\,\rho^a(x')\{\rho^b(y),\rho^c(y')\}$, type $\rho\{\rho,\rho\}$, $(p_1,p_2,p_3)=(x',y,y')$, with
\begin{align}
\mathcal C^{abc}_{\rm I.V}&\equiv \big[U^{ga}(x')-U^{ga}(w')\big]\,f^{gde}\big[U^{dc}(x)\,U^{eb}(z)-U^{eb}(y)\,U^{dc}(y')\big]\,,\notag\\
\mathcal C^{abc}_{\rm I.V}\,\rho^a(x')\{\rho^b(y),\rho^c(y')\}
&=2\,\mathcal C^{abc}_{\rm I.V}\,\{\rho^a(x')\rho^b(y)\rho^c(y')\}_{\rm sym}\notag\\
&\quad+i\,\mathcal C^{abc}_{\rm I.V}\,\Big[f^{abr}\,\delta^{(2)}(x'-y)\,\{\rho^{r}(y)\rho^{c}(y')\}_{\rm sym}\notag\\
&\qquad+f^{acr}\,\delta^{(2)}(x'-y')\,\{\rho^{r}(y')\rho^{b}(y)\}_{\rm sym}\Big]\notag\\
&\quad-\frac16\,\mathcal C^{abc}_{\rm I.V}\big[f^{abr}f^{rcs}+f^{acr}f^{rbs}\big]\notag\\
&\qquad\times\delta^{(2)}(x'-y)\,\delta^{(2)}(x'-y')\,\rho^s(y)\,.\label{eq:app_sym_IV1}
\end{align}
\paragraph{Row II.I, $N_{1a}$, Eq.~\eqref{eq:app_III1a}.} The colour operator, the product of the two brackets of Eq.~\eqref{eq:app_III1a}, is $\mathcal C^{abc}_{1a}\,\{\rho^a(x'),\rho^b(y)\}\rho^c(x)$, type $\{\rho,\rho\}\rho$, $(p_1,p_2,p_3)=(x',y,x)$, with
\begin{align}
\mathcal C^{abc}_{1a}&\equiv U^{ga}(w')\big[f^{kbc}\,U^{gk}(w)-f^{gkl}\,U^{kb}(z)\,U^{lc}(x)\big]\notag\\
&\qquad-U^{ga}(x')\,U^{mh}(z)\,U^{mb}(y)\big[f^{khc}\,U^{gk}(w)-f^{gkl}\,U^{kh}(z)\,U^{lc}(x)\big]\,,\notag\\
\mathcal C^{abc}_{1a}\,\{\rho^a(x'),\rho^b(y)\}\rho^c(x)
&=2\,\mathcal C^{abc}_{1a}\,\{\rho^a(x')\rho^b(y)\rho^c(x)\}_{\rm sym}\notag\\
&\quad+i\,\mathcal C^{abc}_{1a}\,\Big[f^{bcr}\,\delta^{(2)}(y-x)\,\{\rho^{a}(x')\rho^{r}(x)\}_{\rm sym}\notag\\
&\qquad+f^{acr}\,\delta^{(2)}(x'-x)\,\{\rho^{r}(x)\rho^{b}(y)\}_{\rm sym}\Big]\notag\\
&\quad+\frac16\,\mathcal C^{abc}_{1a}\big[f^{bcr}f^{ras}+f^{acr}f^{rbs}\big]\notag\\
&\qquad\times\delta^{(2)}(x'-y)\,\delta^{(2)}(x'-x)\,\rho^s(x)\,.\label{eq:app_sym_III1a}
\end{align}
\paragraph{Row II.I, $N_{1b}$, Eq.~\eqref{eq:app_III1b}.} The colour operator, the product of the two brackets of Eq.~\eqref{eq:app_III1b}, is $\mathcal C^{abc}_{1b}\,\rho^a(x)\{\rho^b(x'),\rho^c(y)\}$, type $\rho\{\rho,\rho\}$, $(p_1,p_2,p_3)=(x,x',y)$, with
\begin{align}
\mathcal C^{abc}_{1b}&\equiv U^{gb}(w)\big[f^{kca}\,U^{gk}(w')-f^{gkl}\,U^{kc}(z)\,U^{la}(x)\big]\notag\\
&\qquad-U^{gb}(x')\,U^{mh}(z)\,U^{mc}(y)\big[f^{kha}\,U^{gk}(w')-f^{gkl}\,U^{kh}(z)\,U^{la}(x)\big]\,,\notag\\
\mathcal C^{abc}_{1b}\,\rho^a(x)\{\rho^b(x'),\rho^c(y)\}
&=2\,\mathcal C^{abc}_{1b}\,\{\rho^a(x)\rho^b(x')\rho^c(y)\}_{\rm sym}\notag\\
&\quad+i\,\mathcal C^{abc}_{1b}\,\Big[f^{abr}\,\delta^{(2)}(x-x')\,\{\rho^{r}(x)\rho^{c}(y)\}_{\rm sym}\notag\\
&\qquad+f^{acr}\,\delta^{(2)}(x-y)\,\{\rho^{r}(x)\rho^{b}(x')\}_{\rm sym}\Big]\notag\\
&\quad-\frac16\,\mathcal C^{abc}_{1b}\big[f^{abr}f^{rcs}+f^{acr}f^{rbs}\big]\notag\\
&\qquad\times\delta^{(2)}(x-x')\,\delta^{(2)}(x-y)\,\rho^s(x)\,.\label{eq:app_sym_III1b}
\end{align}
\paragraph{Row III.II, Eq.~\eqref{eq:app_IIIII}.} The colour operator, the product of the two brackets of Eq.~\eqref{eq:app_IIIII}, is $\mathcal C^{abc}_{\rm III.II}\,\rho^a(y)\rho^b(x')\rho^c(x)$, ordered, $(p_1,p_2,p_3)=(y,x',x)$, with
\begin{align}
\mathcal C^{abc}_{\rm III.II}&\equiv U^{ma}(y)\,U^{gb}(x')\,U^{mh}(z)\big[f^{khc}\,U^{gk}(w)-f^{gkl}\,U^{kh}(z)\,U^{lc}(x)\big]\notag\\
&\qquad-U^{gb}(x')\big[f^{kac}\,U^{gk}(w)-f^{gkl}\,U^{ka}(z)\,U^{lc}(x)\big]\,,\notag\\
\mathcal C^{abc}_{\rm III.II}\,\rho^a(y)\rho^b(x')\rho^c(x)
&=\mathcal C^{abc}_{\rm III.II}\,\{\rho^a(y)\rho^b(x')\rho^c(x)\}_{\rm sym}\notag\\
&\quad+\frac i2\,\mathcal C^{abc}_{\rm III.II}\,\Big[f^{abr}\,\delta^{(2)}(y-x')\,\{\rho^{r}(y)\rho^{c}(x)\}_{\rm sym}\notag\\
&\qquad+f^{bcr}\,\delta^{(2)}(x'-x)\,\{\rho^{a}(y)\rho^{r}(x)\}_{\rm sym}\notag\\
&\qquad+f^{acr}\,\delta^{(2)}(y-x)\,\{\rho^{r}(x)\rho^{b}(x')\}_{\rm sym}\Big]\notag\\
&\quad-\frac16\,\mathcal C^{abc}_{\rm III.II}\big[f^{abr}f^{rcs}-f^{bcr}f^{ras}\big]\notag\\
&\qquad\times\delta^{(2)}(y-x')\,\delta^{(2)}(y-x)\,\rho^s(x)\,.\label{eq:app_sym_IIIII}
\end{align}
\paragraph{Row III.III, first part, structures with $\mathcal L(v)$, $v=y'$ or $z$, Eq.~\eqref{eq:app_III31}.} The colour operators are $f^{c'd'a}\,\mathcal L^{c'd'b}(v)\,\mathcal R_1^{ab}=\mathcal C^{abc}_{\rm III.III}(v)\,\rho^a(x')\{\rho^b(y),\rho^c(x)\}$, type $\rho\{\rho,\rho\}$, $(p_1,p_2,p_3)=(x',y,x)$, with
\begin{align}
\mathcal C^{abc}_{\rm III.III}(v)&\equiv f^{mnk}\big[U^{ma}(v)-U^{ma}(x')\big]\notag\\
&\qquad\times\big[U^{nc}(z)\,U^{kb}(w)-U^{kb}(y)\,U^{nc}(x)\big]\,,\notag\\
\mathcal C^{abc}_{\rm III.III}(v)\,\rho^a(x')\{\rho^b(y),\rho^c(x)\}
&=2\,\mathcal C^{abc}_{\rm III.III}(v)\,\{\rho^a(x')\rho^b(y)\rho^c(x)\}_{\rm sym}\notag\\
&\quad+i\,\mathcal C^{abc}_{\rm III.III}(v)\,\Big[f^{abr}\,\delta^{(2)}(x'-y)\,\{\rho^{r}(y)\rho^{c}(x)\}_{\rm sym}\notag\\
&\qquad+f^{acr}\,\delta^{(2)}(x'-x)\,\{\rho^{r}(x)\rho^{b}(y)\}_{\rm sym}\Big]\notag\\
&\quad-\frac16\,\mathcal C^{abc}_{\rm III.III}(v)\big[f^{abr}f^{rcs}+f^{acr}f^{rbs}\big]\notag\\
&\qquad\times\delta^{(2)}(x'-y)\,\delta^{(2)}(x'-x)\,\rho^s(x)\,.\label{eq:app_sym_IIIIII}
\end{align}
\paragraph{Row III.V, structures with $\mathcal L_{\rm V}(v)$, $v=y'$ or $z$, Eq.~\eqref{eq:app_III5}.} The colour operators are $f^{c'ca}\,\mathcal L_{\rm V}^{c'}(v)\,\mathcal R_{\rm V}^{ac}=\mathcal C^{abc}_{\rm III.V}(v)\,\rho^a(x')\rho^b(x)\rho^c(y)$, ordered, $(p_1,p_2,p_3)=(x',x,y)$, with
\begin{align}
\mathcal C^{abc}_{\rm III.V}(v)&\equiv f^{mkl}\big[U^{ma}(v)-U^{ma}(x')\big]\,U^{kb}(x)\big[U^{lc}(y)-U^{lc}(w)\big]\,,\notag\\
\mathcal C^{abc}_{\rm III.V}(v)\,\rho^a(x')\rho^b(x)\rho^c(y)
&=\mathcal C^{abc}_{\rm III.V}(v)\,\{\rho^a(x')\rho^b(x)\rho^c(y)\}_{\rm sym}\notag\\
&\quad+\frac i2\,\mathcal C^{abc}_{\rm III.V}(v)\,\Big[f^{abr}\,\delta^{(2)}(x'-x)\,\{\rho^{r}(x)\rho^{c}(y)\}_{\rm sym}\notag\\
&\qquad+f^{bcr}\,\delta^{(2)}(x-y)\,\{\rho^{a}(x')\rho^{r}(x)\}_{\rm sym}\notag\\
&\qquad+f^{acr}\,\delta^{(2)}(x'-y)\,\{\rho^{r}(y)\rho^{b}(x)\}_{\rm sym}\Big]\notag\\
&\quad-\frac16\,\mathcal C^{abc}_{\rm III.V}(v)\big[f^{abr}f^{rcs}-f^{bcr}f^{ras}\big]\notag\\
&\qquad\times\delta^{(2)}(x'-x)\,\delta^{(2)}(x'-y)\,\rho^s(x)\,.\label{eq:app_sym_IIIV}
\end{align}
\paragraph{Row III.VI, structures with $\mathcal L(v)$, $v=y'$ or $z$, Eq.~\eqref{eq:app_III5}.} The colour operators are $f^{c'd'a}\,\mathcal L^{c'd'b}(v)\,\mathcal R_{\rm VI}^{ab}=\mathcal C^{abc}_{\rm III.VI}(v)\,\rho^a(x')\rho^b(y)\rho^c(x)$, ordered, $(p_1,p_2,p_3)=(x',y,x)$, with
\begin{align}
\mathcal C^{abc}_{\rm III.VI}(v)&\equiv f^{mnk}\big[U^{ma}(v)-U^{ma}(x')\big]\,U^{kb}(y)\big[U^{nc}(x)-U^{nc}(z)\big]\,,\notag\\
\mathcal C^{abc}_{\rm III.VI}(v)\,\rho^a(x')\rho^b(y)\rho^c(x)
&=\mathcal C^{abc}_{\rm III.VI}(v)\,\{\rho^a(x')\rho^b(y)\rho^c(x)\}_{\rm sym}\notag\\
&\quad+\frac i2\,\mathcal C^{abc}_{\rm III.VI}(v)\,\Big[f^{abr}\,\delta^{(2)}(x'-y)\,\{\rho^{r}(y)\rho^{c}(x)\}_{\rm sym}\notag\\
&\qquad+f^{bcr}\,\delta^{(2)}(y-x)\,\{\rho^{a}(x')\rho^{r}(x)\}_{\rm sym}\notag\\
&\qquad+f^{acr}\,\delta^{(2)}(x'-x)\,\{\rho^{r}(x)\rho^{b}(y)\}_{\rm sym}\Big]\notag\\
&\quad-\frac16\,\mathcal C^{abc}_{\rm III.VI}(v)\big[f^{abr}f^{rcs}-f^{bcr}f^{ras}\big]\notag\\
&\qquad\times\delta^{(2)}(x'-y)\,\delta^{(2)}(x'-x)\,\rho^s(x)\,.\label{eq:app_sym_IIIVI}
\end{align}
\par\endgroup
